\bookchapter{第六十二章：界面之開發——SwiftUI 入門}

\chapterintro{界面者，用戶所見之一切也。SwiftUI 者，蘋果 2019 年所推出聲明式界面之框架也——爾惟「聲明」界面之狀貌，系統自任渲染、適配深色之模式、適配諸般之屏。至 2026 年，SwiftUI 已為新項目絕對之主流。此章，吾儕學「繪畫」。}

\booksection{一、SwiftUI 之思想：聲明式，非命令式}

舊式界面之開發乃「命令式」：先告系統「創按鈕 置於(20,30) 設其色」，一步一步指揮之。SwiftUI 乃「聲明式」：爾徑述「界面之上有一赤色按鈕書『點我』」，如何繪、何時重繪，系統自任之。

\begin{lstlisting}[language=Swift]
import SwiftUI

struct ContentView: View {
    var body: some View {      // body：狀"界面之形"
        VStack {               // 豎而列之
            Text("汝好，世界！")
                .font(.largeTitle)      // 修飾符：大標題
            Button("點我") {
                print("被點擊了")
            }
            .buttonStyle(.borderedProminent)  // 藍色實心按鈕
        }
        .padding()
    }
}
\end{lstlisting}

關鍵之點：視圖者，守 View 協議之結構體也，body 之屬性述界面，.修飾符() 鏈式「上妝」。此即 SwiftUI 全部之基本盤也。

\booksection{二、基礎視圖：文字、按鈕、輸入框、圖}

\begin{center}
\begin{tabularx}{\linewidth}{|X|X|X|}
\hline
\textbf{視圖} & \textbf{何為} & \textbf{示例} \tabularnewline
\hline
Text & 顯示文字 & Text(「你好」) \tabularnewline
Button & 按鈕 & Button(「點我」) \{ 動作 \} \tabularnewline
TextField & 輸入框 & TextField(「提示」, text: \$文本) \tabularnewline
Image & 顯示圖片 & \texttt{Image(systemName: "star.\allowbreak fill")} \tabularnewline
\hline
\end{tabularx}
\end{center}

\begin{center}\small\color{gray}SwiftUI 四大基礎視圖\end{center}

注意 Image(systemName:)——蘋果內置 SF Symbols 圖標之庫（免費圖標數千），「star.fill」此等之名即可徑示一實心之星。此 iOS 開發隱藏之福利：作界面幾乎不必覓圖標之素材。

\booksection{三、佈局三件套：VStack、HStack、ZStack}

\begin{itemize}
  \item VStack：縱而列之（若列隊取餐）；
  \item HStack：橫而列之（若並肩而坐）；
  \item ZStack：疊而置之（若撲克之牌相壘，宜作角標、遮罩）。
\end{itemize}

\begin{lstlisting}[language=Swift]
VStack(alignment: .leading, spacing: 12) {
    HStack {
        Text("賬單")
        Spacer()                 // 彈性之空：推右至盡
        Text("-25.5 元").foregroundColor(.red)
    }
    Divider()                    // 分割線
    Text("備註：午餐")
}
.padding(16)                     // 四周留白
.background(Color(.systemGray6)) // 淺灰底
.cornerRadius(12)                // 圓角

// 修飾符之序有義：先 padding 後 background，背景不含邊距
\end{lstlisting}

\begin{tcbwarn}{按：修飾符次序有講究}
padding(16).background(灰) 與 background(灰).padding(16)，其效迥異——前者先留白而後塗色（背景不含邊距），後者先塗色而後留白（背景兼含邊距）。新手所見怪異之白邊，多源於此。倘有疑，互易其序，預覽立示其答。

\end{tcbwarn}

\booksection{四、狀態與綁定：使界面「活」起來}

靜態界面殊無意味。SwiftUI 之規則：界面恆忠實反映狀態，狀態一變，界面自重繪焉。

\begin{lstlisting}[language=Swift]
struct CounterView: View {
    @State private var count = 0   // @State：此值變則重畫

    var body: some View {
        VStack {
            Text("汝點了 \(count) 次")
            Button("點我 +1") {
                count += 1          // 改狀態，界面自動刷新
            }
        }
    }
}
\end{lstlisting}

輸入框須用「雙向綁定」——\$ 符號將「值與修改之權」一併傳入：

\begin{lstlisting}[language=Swift]
struct InputView: View {
    @State private var text = ""

    var body: some View {
        TextField("輸入若干", text: $text)  // $ 雙向綁定
            .textFieldStyle(.roundedBorder)
        Text("汝輸入了：\(text)")   // 輸入則實時更新
    }
}
\end{lstlisting}

用戶操作

狀態改變（count += 1）

SwiftUI 感知到變化

自動重畫界面（body 重新計算）

\begin{center}\small\color{gray}圖 3：SwiftUI「狀態 界面」單向數據流\end{center}

\booksection{五、預覽、主題與動畫：SwiftUI 之「外掛」}

\begin{itemize}
  \item \#Preview 宏：於檔案之末加 \#Preview \{ ContentView() \}，右側畫布即時渲染——改碼所見即所得；
  \item 深色模式免費：用系統顏色（Color(.systemGray6)、.primary）而非寫死之色，深色模式自能相配，不多書一行；
  \item 動畫一行可成：withAnimation \{ count += 1 \}，或予視圖加 .animation(.default, value: count)——過渡絲滑。
\end{itemize}

\begin{lstlisting}[language=Swift]
struct AnimatedHeartView: View {
    @State private var liked = false

    var body: some View {
        Button {
            withAnimation(.spring(duration: 0.4)) {
                liked.toggle()
            }
        } label: {
            Image(systemName: liked ? "heart.fill" : "heart")
                .font(.system(size: 48))
                .foregroundColor(liked ? .red : .gray)
                .scaleEffect(liked ? 1.3 : 1.0)
        }
    }
}
\end{lstlisting}

\begin{tcbinfo}{小貼士：先功能後動畫}
新手常病：界面未正，而先糾結於動畫。宜先令功能通達，終乃一體加諸動效——動畫者錦上添花也，功能者雪中送炭也。此亦吾儕一貫「簡單至上」之旨。

\end{tcbinfo}

\begin{tcbpractice}{來動手：做一「心情按鈕」}
\booksection{六、常用修飾符速查：一頁在手，界面無憂}

\begin{center}
\begin{tabularx}{\linewidth}{|X|X|}
\hline
\textbf{欲作何事} & \textbf{修飾符} \tabularnewline
\hline
留白 & \texttt{.\allowbreak padding(16)} \tabularnewline
圓角 & \texttt{.\allowbreak cornerRadius(12)} \tabularnewline
背景色 & \texttt{.\allowbreak background(Color(.\allowbreak systemGray6))} \tabularnewline
字體與字號 & \texttt{.\allowbreak font(.\allowbreak title.\allowbreak bold())}/\texttt{.\allowbreak font(.\allowbreak system(size: 20))} \tabularnewline
前景色（文字/圖標） & \texttt{.\allowbreak foregroundColor(.\allowbreak red)} \tabularnewline
撐滿寬度 & \texttt{.\allowbreak frame(maxWidth: .\allowbreak infinity)} \tabularnewline
固定尺寸 & \texttt{.\allowbreak frame(width: 100, height: 44)} \tabularnewline
對齊之法 & \texttt{VStack(alignment: .\allowbreak leading)} \tabularnewline
透明度 & \texttt{.\allowbreak opacity(0.\allowbreak 6)} \tabularnewline
陰影 & \texttt{.\allowbreak shadow(radius: 5)} \tabularnewline
點擊手勢 & \texttt{.\allowbreak onTapGesture \{ .\allowbreak .\allowbreak .\allowbreak  \}} \tabularnewline
鍵盤類型 & \texttt{.\allowbreak keyboardType(.\allowbreak decimalPad)} \tabularnewline
\hline
\end{tabularx}
\end{center}

\begin{center}\small\color{gray}SwiftUI 常用修飾符速查\end{center}

記之不住亦無妨——此頁正為查檢而設。用之最多者，實惟前四：padding、圓角、背景、字體。

改模板之項為心情輪換器：一按鈕，點之則於 開心 難過 生氣 之間循環遞換，文字與背景之色隨之而變。提示：以陣列貯三組（文字+顏色），@State 貯其下標，點時用 (index + 1) \% 3。成則示諸友朋，可得第一聲「哇，汝能作 App 矣」之歎。

\end{tcbpractice}

\begin{tcbsummary}{本章小結}
\begin{itemize}
  \item SwiftUI 為聲明式：述其「所長之形」，系統任描繪重繪之責
  \item 視圖 = 遵 View 協議之結構體；body 述界面；修飾符連綴如「上妝」
  \item 四大基礎視圖：Text / Button / TextField / Image；SF Symbols 圖標庫免費而用之不竭
  \item 佈局三件：VStack 豎列、HStack 橫列、ZStack 疊置；Spacer 者彈性之空也
  \item @State 使界面活起來；\$ 雙向綁定以饋輸入框
  \item \#Preview 所見即所得；系統顏色自配深色模式；withAnimation 一行成動效
\end{itemize}

\end{tcbsummary}
